\chapter{Backend Implementations}

\section{Binary Heap Backend}

The binary heap backend stores entries in a contiguous dynamic array.
Each entry contains the caller-owned item pointer and an optional locator pointer for handled items.

\begin{lstlisting}[style=cstyle, caption={Binary heap entry model.}]
struct binary_heap_handle {
    void *item;
    struct binary_heap_locator *locator;
};
\end{lstlisting}

The array contains no sentinel.
When the heap is non-empty, the minimum item is at index zero.
Insertion appends a new entry and calls \texttt{binary\_heap\_sift\_up()}.
Extract-min removes the root, moves the last entry to the root, and calls \texttt{binary\_heap\_sift\_down()}.

\subsection{Repair After Local Changes}

\texttt{binary\_heap\_repair\_at()} is used after arbitrary removal.
After a removed entry is replaced by the last entry, the replacement may need to move either upward or downward.
The helper checks the parent relation first.
If the replacement compares before its parent, it sifts upward; otherwise it sifts downward.

This keeps arbitrary remove at $O(\log n)$.

\subsection{Capacity Growth}

The binary heap grows its array geometrically.
Capacity starts at eight and doubles when needed.
Size multiplication is checked through the common overflow helper before calling \texttt{realloc}.
After reallocation, locators are refreshed so handled entries still point to the correct array positions.

\section{Fibonacci Heap Backend}

The Fibonacci heap backend stores a circular root list of heap-ordered trees.
Each node stores:

\begin{itemize}
    \item caller-owned item pointer;
    \item optional public handle;
    \item parent pointer;
    \item one child pointer;
    \item left and right sibling pointers;
    \item degree;
    \item mark bit.
\end{itemize}

\subsection{Insertion}

Insertion allocates a node and adds it to the root list.
No consolidation is performed during insertion.
If the new item compares before the current minimum, the minimum pointer is updated.

\subsection{Decrease-Key}

When an item's priority improves, the caller mutates the item and then calls the public decrease-key operation.
The Fibonacci backend checks whether the node now violates heap order with respect to its parent.
If so, it cuts the node from its parent, adds it to the root list, and applies cascading cuts to ancestors that have already lost a child.

This implements the classic mark-bit mechanism used to obtain amortized $O(1)$ decrease-key.

\subsection{Extract-Min and Consolidation}

Extract-min removes the current minimum root, promotes its children to roots, and consolidates the root list by degree.
The implementation snapshots the current roots into a temporary array before linking roots.
This avoids traversing a circular list while mutating it.

The degree table is reused across extract-min calls.
It is sized from the maximum degree and root count seen in the current consolidation.
If allocation for the temporary roots array or degree table fails, the backend preserves correctness by rescanning the root list for the minimum and skipping consolidation for that operation.
This is a deliberate failure-mode choice: correctness is more important than preserving the ideal amortized structure after an allocation failure.

\section{Kaplan Heap Backend}

The Kaplan backend implements the simple Fibonacci heap variant used by this project.
Unlike the classic Fibonacci backend, the steady-state representation is one heap-ordered tree rooted at \texttt{heap->root}.
Children are stored in doubly linked sibling lists.
Each node has a rank and a mark/state bit.

\subsection{Insertion}

Insertion creates a singleton node and links it against the current root.
The better-priority node remains root and the other becomes a child.
This is called a naive link in the implementation.

\subsection{Decrease-Key}

The Kaplan decrease-key repair is different from the Fibonacci backend.
If the decreased node is not the root, the backend performs cascading rank and state changes up the parent chain.
Then it cuts the node once and links it against the root.

This distinction is central to the backend's experimental value.
Both Fibonacci and Kaplan heaps aim at efficient decrease-key behavior, but they repair the structure differently.

\subsection{Extract-Min}

Extract-min removes the root and consolidates its children.
Children are treated as temporary roots.
Roots with equal rank are combined with fair links, incrementing the resulting root's rank.
The remaining roots are then folded into a single tree with naive links.

As with the Fibonacci backend, the rank table is reusable and allocation failures fall back to a correctness-preserving reconstruction.
If the table cannot be allocated, the backend links roots naively.

\subsection{Remove by Handle}

The public API does not expose a backend-independent minus-infinity key.
The Kaplan backend therefore simulates removal structurally.
For a non-root node, it performs the rank repair, detaches the target, makes the old root one of the target's children, installs the target as root, and then reuses extract-min.

This is a good example of why the report should discuss implementation details rather than only public API behavior.
The public operation is simply \texttt{remove(handle)}, but the backend algorithm is specialized.

\section{Backend Comparison}

\begin{table}[H]
    \centering
    \begin{tabular}{llll}
        \toprule
        Backend   & Primary storage     & Target owner   & Main trade-off                \\
        \midrule
        Binary    & Dynamic array       & Stable locator & Locality and simplicity       \\
        Fibonacci & Circular node lists & Node pointer   & Strong amortized bounds       \\
        Kaplan    & Single tree         & Node pointer   & Different decrease-key repair \\
        \bottomrule
    \end{tabular}
    \caption{High-level implementation comparison.}
    \label{tab:backend-comparison}
\end{table}

The binary heap has weaker decrease-key complexity but a simpler memory layout.
The Fibonacci-style heaps have better amortized decrease-key behavior but more metadata and pointer manipulation.
The benchmarks are therefore necessary: asymptotic complexity explains the design space, but implementation details influence measured performance.
